% Options for packages loaded elsewhere
\PassOptionsToPackage{unicode}{hyperref}
\PassOptionsToPackage{hyphens}{url}
\documentclass[
  ignorenonframetext,
]{beamer}
\newif\ifbibliography
\usepackage{pgfpages}
\setbeamertemplate{caption}[numbered]
\setbeamertemplate{caption label separator}{: }
\setbeamercolor{caption name}{fg=normal text.fg}
\beamertemplatenavigationsymbolsempty
% remove section numbering
\setbeamertemplate{part page}{
  \centering
  \begin{beamercolorbox}[sep=16pt,center]{part title}
    \usebeamerfont{part title}\insertpart\par
  \end{beamercolorbox}
}
\setbeamertemplate{section page}{
  \centering
  \begin{beamercolorbox}[sep=12pt,center]{section title}
    \usebeamerfont{section title}\insertsection\par
  \end{beamercolorbox}
}
\setbeamertemplate{subsection page}{
  \centering
  \begin{beamercolorbox}[sep=8pt,center]{subsection title}
    \usebeamerfont{subsection title}\insertsubsection\par
  \end{beamercolorbox}
}
% Prevent slide breaks in the middle of a paragraph
\widowpenalties 1 10000
\raggedbottom
\AtBeginPart{
  \frame{\partpage}
}
\AtBeginSection{
  \ifbibliography
  \else
    \frame{\sectionpage}
  \fi
}
\AtBeginSubsection{
  \frame{\subsectionpage}
}
\usepackage{iftex}
\ifPDFTeX
  \usepackage[T1]{fontenc}
  \usepackage[utf8]{inputenc}
  \usepackage{textcomp} % provide euro and other symbols
\else % if luatex or xetex
  \usepackage{unicode-math} % this also loads fontspec
  \defaultfontfeatures{Scale=MatchLowercase}
  \defaultfontfeatures[\rmfamily]{Ligatures=TeX,Scale=1}
\fi
\usepackage{lmodern}
\ifPDFTeX\else
  % xetex/luatex font selection
\fi
% Use upquote if available, for straight quotes in verbatim environments
\IfFileExists{upquote.sty}{\usepackage{upquote}}{}
\IfFileExists{microtype.sty}{% use microtype if available
  \usepackage[]{microtype}
  \UseMicrotypeSet[protrusion]{basicmath} % disable protrusion for tt fonts
}{}
\makeatletter
\@ifundefined{KOMAClassName}{% if non-KOMA class
  \IfFileExists{parskip.sty}{%
    \usepackage{parskip}
  }{% else
    \setlength{\parindent}{0pt}
    \setlength{\parskip}{6pt plus 2pt minus 1pt}}
}{% if KOMA class
  \KOMAoptions{parskip=half}}
\makeatother
\setlength{\emergencystretch}{3em} % prevent overfull lines
\providecommand{\tightlist}{%
  \setlength{\itemsep}{0pt}\setlength{\parskip}{0pt}}
% Auto-generated by infrastructure/rendering/_pdf_latex_helpers.py::extract_math_font_preamble.
% Minimal header passed to Pandoc via -H so Beamer slides pick up the same math font
% as the combined-PDF path. Adjust the manuscript's preamble.md to change the font globally.
\usepackage{unicode-math}
\setmathfont{latinmodern-math.otf}

% Slide layout defaults for warning-clean scientific decks.
\usepackage{etoolbox}
\IfFileExists{xurl.sty}{\usepackage{xurl}}{}
\IfFileExists{seqsplit.sty}{\usepackage{seqsplit}}{\newcommand{\seqsplit}[1]{#1}}
\protected\def\breakseq#1{\seqsplit{#1}}
\protected\def\breaktt#1{\begingroup\ttfamily\seqsplit{#1}\endgroup}
\setlength{\emergencystretch}{6em}
\tolerance=5000
\hbadness=10000
\hfuzz=1pt
\setlength{\tabcolsep}{2pt}
\AtBeginEnvironment{longtable}{\tiny\renewcommand{\arraystretch}{0.86}\setlength{\tabcolsep}{1pt}}
\AtBeginEnvironment{tabular}{\tiny\renewcommand{\arraystretch}{0.86}\setlength{\tabcolsep}{1pt}}

% Natbib and cross-reference fallbacks — slides don't load natbib
% or cleveref, but combined-PDF manuscript prose may emit these
% commands. The fallback renders citations as a bracketed key list
% and cross-references as detokenized labels so slides don't fail on
% undefined control sequences or raw underscores. \providecommand is
% a no-op if packages load later.
\providecommand{\citep}[1]{[#1]}
\providecommand{\citet}[1]{#1}
\providecommand{\citealp}[1]{#1}
\providecommand{\citeauthor}[1]{#1}
\providecommand{\citeyear}[1]{#1}
\providecommand{\cref}[1]{\texttt{\detokenize{#1}}}
\providecommand{\Cref}[1]{\texttt{\detokenize{#1}}}
\usepackage{bookmark}
\IfFileExists{xurl.sty}{\usepackage{xurl}}{} % add URL line breaks if available
\urlstyle{same}
% fallback for those not using the hyperref driver hyperxmp:
\makeatletter
\@ifundefined{xmpquote}{\newcommand{\xmpquote}[1]{#1}}{}
\makeatother
\hypersetup{
  hidelinks,
  pdfcreator={LaTeX via pandoc}}

\author{\texorpdfstring{}{}}
\date{}

\begin{document}

\section{Methodology}\label{sec:methodology}

\begin{frame}[fragile,allowframebreaks]{Methodology}
This section describes the implementation methodology, explicitly
detailing how the optimization algorithms are constructed, validated,
and analyzed using the Generalized Research Template's
\texttt{infrastructure} and \texttt{tests} ecosystems.
\end{frame}

\begin{frame}[fragile,allowframebreaks]{Algorithm Implementation}
\protect\phantomsection\label{algorithm-implementation}
\begin{block}{Gradient Descent Algorithm}
\protect\phantomsection\label{gradient-descent-algorithm}
The core algorithm implements the iterative procedure for unconstrained
optimization. The
\href{https://github.com/docxology/template/blob/main/projects/templates/template_code_project/src/optimizer.py}{\texttt{optimizer}
module} uses the standard-library \texttt{logging} logger for optional
verbose diagnostics; the
\href{https://github.com/docxology/template/blob/main/projects/templates/template_code_project/scripts/optimization_analysis.py}{analysis
orchestrator} uses
\breaktt{infrastructure.core.logging.utils.get\_logger}. Tests run under
the hermetic boundaries defined in the
\href{https://github.com/docxology/template/blob/main/projects/templates/template_code_project/tests/conftest.py}{test
configuration}.

\textbf{Algorithm --- Gradient Descent (implemented in the
\href{https://github.com/docxology/template/blob/main/projects/templates/template_code_project/src/optimizer.py\#L87-L173}{optimizer
module})}

\begin{quote}
\textbf{Input:} Initial point \(x_0\), step size \(\alpha\), tolerance
\(\epsilon\), max iterations \(N_{\max}\)

\begin{enumerate}
\tightlist
\item
  Initialize \(k \leftarrow 0\)
\item
  \textbf{While} \(k < N_{\max}\) \textbf{do:}

  \begin{itemize}
  \tightlist
  \item
    Compute gradient \(g_k = \nabla f(x_k)\)
  \item
    \textbf{If} \(\|g_k\|_2 < \epsilon\) \textbf{then return} \(x_k\)
    \emph{(converged)}
  \item
    Update: \(x_{k+1} \leftarrow x_k - \alpha \cdot g_k\)
  \item
    \(k \leftarrow k + 1\)
  \end{itemize}
\item
  \textbf{Return} \(x_k\) \emph{(max iterations reached)}
\end{enumerate}
\end{quote}
\end{block}
\end{frame}

\begin{frame}[fragile,allowframebreaks]{Infrastructure Integration}
\protect\phantomsection\label{infrastructure-integration}
The methodology explicitly bridges theoretical mathematics with
production-grade validation through the
\breaktt{infrastructure.scientific} module.

\begin{block}{Numerical Stability Analysis}
\protect\phantomsection\label{numerical-stability-analysis}
Rather than writing ad-hoc validation code, the project imports
\breaktt{infrastructure.scientific.stability.check\_numerical\_stability}.
This utility subjects the objective function to a barrage of extreme
inputs (NaN, Inf, \(\pm 10^{10}\)) to calculate a formalized stability
score. If this score degrades, the
\href{https://github.com/docxology/template/blob/main/scripts/02_run_analysis.py}{analysis
orchestrator} execution deliberately aborts, ensuring the methodology
cannot enter unrecoverable states.
\end{block}

\begin{block}{Performance Benchmarking}
\protect\phantomsection\label{performance-benchmarking}
Computational complexity is evaluated not just theoretically, but
empirically via
\href{https://github.com/docxology/template/blob/main/infrastructure/scientific/benchmarking.py}{\breaktt{infrastructure.scientific.benchmarking.benchmark\_function}}.
This module captures high-resolution execution timings and memory
footprints across dimensionality sweeps, guaranteeing that the \(O(n)\)
space-time complexity predictions hold true on the host architecture.
\end{block}
\end{frame}

\begin{frame}[allowframebreaks]{Convergence Analysis}
\protect\phantomsection\label{convergence-analysis}
For quadratic functions \(f(x) = \frac{1}{2}x^T A x - b^T x\) where
\(A\) is positive definite, the convergence factor becomes
\breakseq{[@bertsekas1999nonlinear]}:

\begin{equation}
\label{eq:convergence_factor}
\rho = \frac{|\lambda_{\max} - \alpha\lambda_{\min}|}{|\lambda_{\min} + \alpha\lambda_{\max}|}
\end{equation}

Optimal convergence occurs when
\(\alpha = \frac{2}{\lambda_{\min} + \lambda_{\max}}\), yielding
\(\rho = \frac{\kappa - 1}{\kappa + 1}\).
\end{frame}

\begin{frame}[fragile,allowframebreaks]{Experimental Setup}
\protect\phantomsection\label{experimental-setup}
\begin{block}{Step Size Analysis}
\protect\phantomsection\label{step-size-analysis}
Step sizes are not chosen ad hoc in the manuscript: they are read from
\breaktt{experiment.step\_sizes} in \breaktt{manuscript/config.yaml} and
passed through \breaktt{run\_convergence\_experiment()} in
\texttt{src/analysis/} (entry:
\breaktt{scripts/optimization\_analysis.py}). The active grid for this
build is:

\begin{itemize}
\tightlist
\item
  \(\alpha = 0.01\) (conservative)
\item
  \(\alpha = 0.1\) (conservative)
\item
  \(\alpha = 0.5\) (near-optimal)
\item
  \(\alpha = 1.0\) (near-optimal)
\item
  \(\alpha = 1.5\) (aggressive)
\item
  \(\alpha = 2.5\) (divergent (expected unstable for H = I))
\end{itemize}

Labels follow the same agency taxonomy used for plot colours
(\breaktt{\_agency\_category} in the analysis script):
\textbf{conservative} (small \(\alpha\), slow but stable on \(H=I\)),
\textbf{near-optimal} (including \(\alpha=1\) where the method reaches
\(x^\ast\) in one step for this quadratic), \textbf{aggressive}
(\(1 < \alpha < 2\), still linearly contracting in \(|x-x^\ast|\) but
oscillatory in sign), and \textbf{divergent} (\(|1-\alpha| \geq 1\)).
\end{block}

\begin{block}{Zero-Mock Testing Methodology}
\protect\phantomsection\label{zero-mock-testing-methodology}
The most critical aspect of the project's methodology is its validation
framework. The project is governed by a strict Zero-Mock testing policy,
evaluated actively by executing
\breaktt{uv\ run\ pytest\ projects/templates/template\_code\_project/tests/}
during the infrastructure build phase.

\begin{enumerate}
\tightlist
\item
  \textbf{Project tests}:
  \href{../tests/test_optimizer.py}{\breaktt{projects/templates/template\_code\_project/tests/test\_optimizer.py}}
  exercises \breaktt{src/optimizer.py} (typical, edge, boundary, and
  pathological inputs including NaN/Inf and zero gradients) and, when
  infrastructure imports succeed, call into
  \breaktt{optimization\_analysis.py} helpers---without mocks. Suite
  size:
  \href{../../../../docs/_generated/COUNTS.md}{\breaktt{docs/\_generated/COUNTS.md}}.
\item
  \textbf{Infrastructure validation}: The repository-level
  \breaktt{tests/infra\_tests/} suite validates shared template modules
  (e.g.~pipeline and discovery helpers) independently of this project's
  manuscript.
\item
  \textbf{Coverage Gates}: The
  \href{https://github.com/docxology/template/blob/main/.github/workflows/ci.yml}{GitHub
  Actions CI workflow} enforces a mandatory ≥90\% statement coverage
  gate on \breaktt{projects/templates/template\_code\_project/src/} prior
  to treating the project as build-green.
\end{enumerate}
\end{block}

\begin{block}{Stopping rule and reporting}
\protect\phantomsection\label{stopping-rule-and-reporting}
\breaktt{gradient\_descent()} terminates when \(\|\nabla f(x_k)\|\) falls
below \breaktt{experiment.tolerance} or when \(k\) reaches
\breaktt{experiment.max\_iterations}. The boolean \texttt{converged} in
exported CSV rows distinguishes these outcomes. Downstream,
\breaktt{scripts/z\_generate\_manuscript\_variables.py} aggregates the
CSV into \texttt{RESULT\_*} placeholders so tables and prose cannot
drift from the last analysis run.
\end{block}

\begin{block}{Figure generation contract}
\protect\phantomsection\label{figure-generation-contract}
Each figure in \texttt{03\_results.md} maps to a generator in
\texttt{src/figures/} (\breaktt{generate\_convergence\_plot},
\breaktt{generate\_step\_size\_sensitivity\_plot},
\breaktt{generate\_convergence\_rate\_plot},
\breaktt{generate\_complexity\_visualization},
\breaktt{generate\_stability\_visualization},
\breaktt{generate\_benchmark\_visualization}), orchestrated by
\texttt{src/analysis/} / \breaktt{scripts/optimization\_analysis.py}.
Captions in the markdown intentionally name the function and the key
parameters (tolerance lines, grids, dimensions) so reviewers can
navigate from PDF to code without inferring hidden defaults.
\end{block}
\end{frame}

\begin{frame}[fragile,allowframebreaks]{Analysis Pipeline \& LaTeX
Integration}
\protect\phantomsection\label{analysis-pipeline-latex-integration}
The automated analysis script leverages
\href{https://github.com/docxology/template/blob/main/infrastructure/core/progress.py}{\breaktt{infrastructure.core.progress}}
(\texttt{ProgressBar}, \breaktt{SubStageProgress}) to orchestrate
experiments, collect convergence trajectories, and generate
publication-quality visualizations seamlessly.

The research template supports advanced LaTeX customization through the
\texttt{preamble.md} configuration. This is ingested directly by
\breaktt{infrastructure.rendering.latex\_utils.py} and
\texttt{pdf\_renderer.py}, automatically linking compiled PGF plots and
BibTeX citations. This automated approach ensures an unbreakable chain
of custody from raw algorithmic execution to the final rendered
manuscript.
\end{frame}

\end{document}
